\documentclass[10pt,a4paper]{article}

\usepackage[margin=22mm]{geometry}
\usepackage{fontspec}
\usepackage{booktabs}
\usepackage{tabularx}
\usepackage{enumitem}
\usepackage{xcolor}
\usepackage{hyperref}
\usepackage{microtype}

\setmainfont{TeX Gyre Termes}
\setsansfont{TeX Gyre Heros}
\hypersetup{colorlinks=true,urlcolor=blue,linkcolor=black}
\setlist{nosep,leftmargin=1.5em}
\setlength{\parindent}{0pt}
\setlength{\parskip}{4pt}
\setlength{\emergencystretch}{3em}
\newcommand{\code}[1]{\texttt{\small #1}}

\title{\textbf{Multi-Camera Multi-Vehicle Tracking Dataset and Algorithm Brief}\\
\large NDNSF-UAV integration note}
\author{NDNSF-UAV project}
\date{September 2026}

\begin{document}
\maketitle

\textbf{Purpose.} This note describes the public ``Multi-Camera Multi-Vehicle
Tracking System (UAV Dataset)'' and its reference tracking pipeline, and maps
them to a bounded NDNSF-UAV experiment. The dataset is useful for a small,
reproducible multi-provider prototype; it is not by itself a complete
multi-UAV scientific benchmark for recognition or geolocation.

\section{Dataset overview and provenance}

The dataset is distributed with the project
\href{https://huggingface.co/datasets/jye9/Multi-Camera-Multi-Vehicle-Tracking-System}{Hugging Face dataset card}.
The accompanying implementation is available at
\href{https://github.com/JYe9/multi-camera-multi-vehicle-tracking-system}{the project repository}.
The repository describes three synchronized one-minute UAV clips:
\code{uav1\_1min.mp4}, \code{uav2\_1min.mp4}, and \code{uav3\_1min.mp4}.
The views observe overlapping urban regions and are intended to maintain a
consistent identity for the same vehicle across cameras. The CSV output uses
frame-level fields such as bounding box, confidence, local tracker ID,
\code{global\_id}, \code{camera\_id}, speed, and status.

\begin{table}[h]
\centering
\small
\begin{tabularx}{\linewidth}{@{}lX@{}}
\toprule
Property & Dataset-level interpretation \\
\midrule
Views & Three UAV video streams with synchronized frame indices \\
Target & Vehicles in urban scenes; the example highlights a red vehicle with a persistent global ID \\
Labels & Per-frame boxes and tracking metadata; local and cross-camera IDs are both represented \\
Output & Annotated MP4 files and per-camera CSV files \\
Algorithm context & Ultralytics YOLO detection, Supervision/ByteTrack local tracking, and overlap-region cross-camera matching \\
Scale & Small pilot resource: three one-minute videos rather than a large, diverse flight campaign \\
License boundary & Check the dataset card and project license before redistribution; keep raw videos outside Git \\
\bottomrule
\end{tabularx}
\end{table}

The data should therefore be called a \emph{synchronized multi-camera UAV
tracking pilot}. It should not be described as a broad recognition benchmark:
the release does not provide the target-class diversity, calibrated camera
geometry, or flight metadata needed for a general claim about 3-D localization
or improved recognition accuracy. It is also different from ordinary UAV MOT:
multi-object tracking within one video does not establish identity association
across independent views.

\clearpage
\section{Reference algorithm}

The reference code is an engineering baseline rather than a new tracking
algorithm. Its processing chain is:

\begin{enumerate}
  \item Load one YOLO detector and initialize one local tracker per UAV stream.
  \item Process corresponding frame indices in parallel. Each stream performs
        confidence filtering, non-maximum suppression, and local ByteTrack
        association.
  \item Convert local detections into anchor points and estimate speed and
        direction using camera-specific pixel-to-meter scale factors.
  \item Use configured field-of-view and overlap regions to match compatible
        local tracks across cameras and assign a persistent \code{global\_id}.
  \item Append detection and association records to CSV and render boxes,
        labels, traces, regions, counters, and time codes to annotated video.
\end{enumerate}

The important distinction is between three identities: the detector's object
candidate, a camera-local \code{tracker\_id}, and the cross-camera
\code{global\_id}. A global ID is an application-level association result; it
is not automatically an NDN producer identity and it is not a cryptographic
proof that two objects are physically the same vehicle. The matching logic is
based on overlap regions and configured spatial/appearance cues, so its speed
and direction estimates require scene-specific calibration.

For an NDNSF experiment, each camera stream should be treated as a separate
logical Provider. The application should publish a bounded frame or temporal
chunk as Named Data, while the tracking service consumes references to those
objects. The full video should not be placed in a Request or Selection
message.

\clearpage
\section{NDNSF-UAV integration and evaluation}

\subsection*{3.1 Data-plane mapping}

The three streams can be mapped to a minimal collaboration plan:

\begin{verbatim}
/uav/A  -> ViewData-A  --+
/uav/B  -> ViewData-B  ---+--> TrackingProvider -> User
/uav/C  -> ViewData-C  --+
\end{verbatim}

The protected discovery descriptor may contain the target region, deadline,
minimum number of views, required vehicle-tracking capability, and resource
constraints. After ACK collection and Selection, the plan assigns capture
roles to the selected UAV Providers and a tracking role to the analysis
Provider. Each View Data object should carry at least:

\begin{itemize}
  \item a request/mission session and attempt identifier;
  \item producer identity, camera/UAV identity, sequence and frame index;
  \item capture timestamp, target or track identifier, and media type;
  \item exact Data name, content digest, bounding-box metadata, and optional
        pose/GPS information.
\end{itemize}

The tracker must accept a View Data object only when its producer, name, digest,
capture window, and dependency edge agree with the selected plan. Directly
retrieved and cached objects must pass the same checks. A missing Provider,
wrong producer, undeclared dependency, duplicate object, or stale attempt is a
protocol rejection—not a computer-vision false negative.

\subsection*{3.2 Recommended experiment stages}

\begin{enumerate}
  \item \textbf{Pipeline smoke test.} Use the three released videos to exercise
        three ACKs, three selected capture roles, Named Data publication, and
        one tracking Provider. Verify direct retrieval, cache retrieval, and
        result publication.
  \item \textbf{Failure and admission tests.} Remove one stream, delay one
        stream, substitute a wrong producer, replay an old frame, or submit an
        undeclared dependency. Record the first rejection stage and the final
        request outcome.
  \item \textbf{Paired tracking comparison.} Compare one stream against two
        and three streams over the same time windows. Keep detector weights,
        thresholds, frame sampling, and preprocessing fixed.
  \item \textbf{Extended application study.} Only after obtaining more scenes
        and target identities should the project report general multi-view
        tracking or recognition improvements. The current three one-minute
        release is a pilot, not sufficient evidence for such a claim.
\end{enumerate}

\subsection*{3.3 Metrics and limitations}

Application metrics should include IDF1, HOTA or MOTA, identity switches,
cross-view association accuracy, track continuity, and recovery after
occlusion. NDNSF metrics should include positive ACK count, selected Provider
count, accepted-view count, completion within deadline, end-to-end latency,
message/byte overhead, cache-hit latency, and rejection counts for wrong
producer, wrong dependency, and stale attempt.

The dataset is especially useful for validating the \emph{data-plane contract}
of NDNSF-UAV: selected Providers publish named objects and the analysis role
advances only on authorized, correctly correlated objects. It does not by
itself validate UAV flight control, GPS localization accuracy, cryptographic
authorization strength, or that adding views always improves recognition.
Those claims require calibrated metadata, more independent scenes, target-level
splits, and a separately controlled application study.

\section*{Sources}

\begin{itemize}
  \item Dataset card: \href{https://huggingface.co/datasets/jye9/Multi-Camera-Multi-Vehicle-Tracking-System/blob/main/README.md}{Hugging Face README}
  \item Algorithm repository: \url{https://github.com/JYe9/multi-camera-multi-vehicle-tracking-system}
  \item Project paper record: \url{https://arxiv.org/abs/2605.15779}
\end{itemize}

\end{document}
